\chapter{Montgomery Curves}
\label{chap:montgomeryCurves}

In our voting scheme, we use \gls{eeg} encryption for the ballots.
All operation in the encryption of a value in \gls{eeg} are performed in a finite cyclic group $\mathbb{G}$.
For a secure encryption, the \gls{ddh} assumption needs to hold on $\mathbb{G}$.
Elliptic curves of a sufficiently large order $q$ are finity cyclic groups on which the \gls{ddh} assumption is believed to hold.
Moreover, Elliptic curves are benefitial compared to other cyclic groups because they need a smaller order to produce similar levels of security \cite{noauthor_case_2009}.

The most general form of Elliptic curves are so-called Short Weierstrass curves \cite{least_authority_moonmath_2024}.
Montgomery Curves are a subclass of Short Weierstrass curves \cite{montgomery_speeding_nodate}.
We can transform any Montgomery curve into an isomorphic Short Weierstrass curve but not every Short Weierstrass curves has an isomorphic Montgomery counterpart.
Therefore, Montgomery curves are more limited than Short Weierstrass curves \cite{least_authority_moonmath_2024}.

However, as established previously, we will later use an Elliptic curve as the group $(\mathbb{G}, \odot_{\mathbb{G}})$ for \gls{eeg} encryption.
In \gls{eeg} we need to compute exponentiations of group elements, i.e., $g^m:=g\underbrace{\odot_{\mathbb{G}}\dots\odot_{\mathbb{G}}}_{m \text{ times}} g$ very often.
Montgomery curves allow us to compute $g^m$ efficiently using the Montgomery Ladder algorithm.

Fundamentally, Elliptic curves are sets of points that satisfy certain equations \cite{least_authority_moonmath_2024}. 
Here, a curve point consists of multiple elements from a so-called base field $\mathbb{F}_p$ of order $p$.
For better distinction, we call the field $\mathbb{F}_q$ with elliptic curve order $q$ the scalar field of the elliptic curve.

Our introduction of elliptic curves largely parallels the one given in chapter 5 of \cite{least_authority_moonmath_2024}.
However, as we exclusively use Montgomery Curves, we introduce them directly and do not introduce Short Weierstrass Curves seperately.

\section{Affine Montgomery Form}
As in \cite{least_authority_moonmath_2024} we introduce elliptic curves first in their affine representation and then in their projective representation.
Later on, we use both representations in the voting scheme.

\begin{definition}[Affine Montgomery Curve]
    Let $(\mathbb{F}, +, \cdot)$ be a prime field of order $p>3$.
    Let $A, B\in\mathbb{F}$ be two field elements with $B\neq 0$ and $A^2\neq4 \mod p$.
    Then, the corresponding Affine Montgomery curve is the group $(M_{A,B}(\mathbb{F}), \odot_M^{Aff})$, with:
    \begin{itemize}
        \item The set of points $M_{A,B}(\mathbb{F})$:
        \begin{align*}
            M_{A,B}(\mathbb{F})=\{(x,y)\in\mathbb{F}\times\mathbb{F}|B\cdot y^2=x^3+A\cdot x^2+x\}\cup \{\mathcal{O}\}
        \end{align*}
        Here, $\mathcal{O}$ is the so-called \glqq point at infinity\grqq.
        \item The group law $\odot_M^{Aff}$:\\
        Let $P_1, P_2\in M_{A,B}(\mathbb{F})$.
        Then:
        \begin{align*}
            P_1\odot_M^{Aff} P_2 =\begin{cases}
                P_1, &\quad\text{if } P_2=\mathcal{O}\\
                P_2, &\quad\text{if } P_1=\mathcal{O}\\
                \mathcal{O}, &\quad\text{if } P_1=(x, y) \text{ and } P_2=(x, -y)\\
                TangentRule(x,y), &\quad\text{if } P_1=P_2=(x,y) \text{ and } y\neq 0\\
                ChordRule(x_1, y_1, x_2, y_2), &\quad\text{otherwise, with } P_1=(x_1, y_1), P_2=(x_2, y_2)
                %\left(\left(\frac{3x^2+a}{2y}\right)^2-2x, \left(\frac{3x^2+a}{2y}\right)\cdot\left(x-\left(\left(\frac{3x^2+a}{2y}\right)^2-2x\right)\right)-y\right), &\quad\text{if } P_1=P_2=(x,y) \text{ and } y\neq 0\\
                %\left(\left(\frac{y_2-y_1}{x_2-x_2}\right)^2-x_1-x_2, \left(\frac{y_2-y_1}{x_2-x_1}\right)\cdot\left(x_1-\left(\left(\frac{y_2-y_1}{x_2-x_2}\right)^2-x_1-x_2\right)\right)-y_1\right), &\quad\text{otherwise, with } P_1=(x_1, y_1), P_2=(x_2, y_2)
            \end{cases}
        \end{align*}
        $TangentRule$ and $ChordRule$ are defined in \ref{alg:tangentRuleAffineMontgomery} and \ref{alg:chordRuleAffineMontgomery}.
    \end{itemize}
    The neutral element of $M_{A,B}(\mathbb{F})$ is $\mathcal{O}$.
    For $P\in M_{A,B}(\mathbb{F})$, the inverse element $P^{-1}$ is defined as $P^{-1}:=\begin{cases}
        \mathcal{O}, &\quad\text{if } P=\mathcal{O}\\
        (x, -y), &\quad\text{otherwise, with } P=(x,y)
    \end{cases}$.
\end{definition}

\begin{Algorithmus} %Use the environment only if you want to place the algorithm similar to graphics from TeX
    \caption{Affine Tangent Rule used to compute $P\odot_M^{Aff} P$ where $P=(x,y)\in M_{A,B}(\mathbb{F})\setminus\{\mathcal{O}\}$ and $y\neq 0$.}
    \label{alg:tangentRuleAffineMontgomery}
    \begin{algorithmic}
    \Procedure{$TangentRule$}{$x,y$}
        \State $tmp := \frac{3x^2+2Ax+1}{2By}$
        \State $x' := tmp^2\cdot B-2x-A$
        \State $y':= tmp\cdot (x-x')-y$
        \State \Return $(x', y')$
    \EndProcedure
    \end{algorithmic}
\end{Algorithmus}

\begin{Algorithmus} %Use the environment only if you want to place the algorithm similar to graphics from TeX
    \caption{Affine Chord Rule used to compute $P_1\odot_M^{Aff} P_2$ where $P_1=(x_1,y_1), P_2=(x_2, y_2)\in M_{A,B}(\mathbb{F})\setminus\{\mathcal{O}\}$ and $x_1\neq x_2$.}
    \label{alg:chordRuleAffineMontgomery}
    \begin{algorithmic}
    \Procedure{$ChordRule$}{$x_1,y_1, x_2, y_2$}
        \State $tmp := \frac{y_2-y_1}{x_2-x_1}$
        \State $x' := tmp^2\cdot B-x_1-x_2-A$
        \State $y':= tmp\cdot (x_1-x')-y_1$
        \State \Return $(x', y')$
    \EndProcedure
    \end{algorithmic}
\end{Algorithmus}

%\subsection{Affine and Projective Coordinates}
%To avoid the special point at infinity $\mathcal{O}$ there exists the projective representation of elliptic curves based on projective planes.

\section{Projective Montgomery Form}
To avoid the special point at infinity $\mathcal{O}$ there exists the projective representation of elliptic curves based on projective planes.

\begin{definition}[Projective Point and Projective Plane]
    Let $(\mathbb{F}, +, \cdot)$ be a field. 
    We define $\mathbb{F}^*$ as the set of invertible elements in $\mathbb{F}$.
    Let $(X, Y, Z)=\in\mathbb{F}^3$. 
    Then, the projective point corresponding to $(X, Y, Z)$ is defined as
    \begin{align*}
        [X:Y:Z] :=\{(k\cdot X, k\cdot Y, k\cdot Z)|k\in\mathbb{F}^*\}\text{.}
    \end{align*}
    The projective plane of $\mathbb{F}$ is defined as
    \begin{align*}
        \mathbb{F}P^2:=\{[X:Y:Z]|(X,Y,Z)\in\mathbb{F}^3 \text{ and } (X, Y, Z)\neq (0,0,0)\}\text{.}
    \end{align*}
\end{definition}

Using projective planes, we can define Montgomery Curves in their projective representation as groups of projective points.

\begin{definition}[Projective Montgomery Curve]
    Let $(\mathbb{F}, +, \cdot)$ be a prime field of order $p>3$. 
    Let $A, B\in\mathbb{F}$ be two field elements with $B\neq 0$ and $A^2\neq4 \mod p$.
    Let $\mathbb{F}P^2$ be the projective plane of $\mathbb{F}$.
    Then, the corresponding Projective Montgomery curve is the group $(M_{A,B}(\mathbb{F}P^2), \odot_M^{Pro})$, with:
    \begin{itemize}
        \item The set of points $M_{A,B}(\mathbb{F}P^2)$:
        \begin{align*}
            M_{A,B}(\mathbb{F}P^2)=\{[X:Y:Z]\in\mathbb{F}P^2|B\cdot Y^2\cdot Z=X^3+A\cdot X^2\cdot Z+X\cdot Z^2\}
        \end{align*}
        Here, $[0:1:0]$ corresponds to the point at infinity $\mathcal{O}$ of the affine form.
        \item The group law $\odot_M^{Pro}$:\\
        Let $P_1=[X_1:Y_1:Z_1], P_2=[X_2:Y_2:Z_2]\in M_{A,B}(\mathbb{F}P^2)$. 
        To compute $P_1\odot_M^{Pro} P_2$ we transform $P_1, P_2$ to corresponding points on the corresponding Affine Montgomery Curve $M_{A,B}(\mathbb{F})$, use the affine group law $\odot_{M}^{Aff}$ and transform the result back to $M_{A,B}(\mathbb{F}P^2)$.
        Formally:
        \begin{align*}
            P_1\odot_M^{Pro} P_2 = I(I^{-1}(P_1)\odot_M^{Aff}I^{-1}(P_2))
        \end{align*}
        % $TangentRule$ and $ChordRule$ are defined in \ref{alg:tangentRuleProjectiveMontgomery} and \ref{alg:chordRuleProjectiveMontgomery}.
        Here, $I$ converts points from affine representation to projective representation, and $I^{-1}$ executes the reverse conversion.
        $I$ and $I^{-1}$ are introduced in \cref{sec:coordinateTranformation}.
    \end{itemize}
    The neutral element of $M_{A,B}(\mathbb{F}P^2)$ is $[0:1:0]$.
    For $P=[X:Y:Z]\in M_{A,B}(\mathbb{F}P)$, the inverse element $P^{-1}$ is defined as $P^{-1}:=[X:-Y:Z]$.
\end{definition}

\section{Transforming between Projective and Affine Coordinates}
\label{sec:coordinateTranformation}
As mentioned earlier, we will use both the projective and affine representation of elliptic curves in our voting schemes.
Therefore, we need to be able to convert from an affine elliptic curve $M_{A,B}(\mathbb{F})$ to a corresponding projective elliptic curve $M_{A,B}(\mathbb{F}P^2)$ and vice versa.

\begin{definition}[Coordinate Transformation]
    Let $(\mathbb{F}, +, \cdot)$ be a finite field of order $q$ with $q>3$. 
    Let $a,b\in \mathbb{F}$ with $4a^2 + 27b^2\neq 0$. 
    Let $\mathbb{F}P^2$ be the projective plane of $\mathbb{F}$.
    Then, the bijective function $I:M_{A,B}(\mathbb{F})\to M_{A,B}(\mathbb{F}P^2)$ converts points from affine representation to projective representation.
    Consequently, $I^{-1}:M_{A,B}(\mathbb{F}P^2)\to M_{A,B}(\mathbb{F})$ converts points from projective representation to affine representation.
    Here, we define:
    \begin{align*}
        I(P) &:= \begin{cases}
            [x:y:1], \quad\text{if } P\neq\mathcal{O}, \text{ with } P=(x,y)\\
            [0:1:0], \quad\text{otherwise}
        \end{cases}\\
        I^{-1}([X:Y:Z]) &:= \begin{cases}
            \left(\frac{X}{Z}, \frac{Y}{Z}\right), \quad\text{if } Z\neq 0\\
            \mathcal{O}, \quad\text{otherwise}
        \end{cases}
    \end{align*}
\end{definition}

%Note that this conversion works the same for Short Weierstrass elliptic curves and Montgomery Curves which are introduced next.

% \begin{Algorithmus} %Use the environment only if you want to place the algorithm similar to graphics from TeX
%     \caption{Projective Tangent Rule used to compute $P\odot_M^{Pro} P$ where $P=[X:Y:Z]\in M_{A,B}(\mathbb{F})\setminus\{[0:1:0]\}$ and $Y\neq 0$.}
%     \label{alg:tangentRuleProjectiveMontgomery}
%     \begin{algorithmic}
%     \Procedure{$TangentRule$}{$[X:Y:Z]$}
%         \State \todo{Write Alg.}
%         \State \Return $([X':Y':Z'])$
%     \EndProcedure
%     \end{algorithmic}
% \end{Algorithmus}
% 
% \begin{Algorithmus} %Use the environment only if you want to place the algorithm similar to graphics from TeX
%     \caption{Projective Chord Rule used to compute $P_1\odot_M^{Pro} P_2$ where $P_1=[X_1:Y_1:Z_1], P_2=[X_2:Y_2:Z_2]\in M_{A,B}(\mathbb{F}P^2)\setminus\{[0:1:0]\}$ and $X_1\neq X_2$.}
%     \label{alg:chordRuleProjectiveMontgomery}
%     \begin{algorithmic}
%     \Procedure{$ChordRule$}{$[X_1:Y_1:Z_1], [X_2:Y_2:Z_2]$}
%         \State \todo{Write Alg.}
%         \State \Return $([X':Y':Z'])$
%     \EndProcedure
%     \end{algorithmic}
% \end{Algorithmus}



% \begin{Algorithmus} %Use the environment only if you want to place the algorithm similar to graphics from TeX
%     \caption{Projective Tangent Rule used to compute $P\odot_M^{Pro} P$ where $P=[X:Y:Z]\in M_{A,B}(\mathbb{F})\setminus\{[0:1:0]\}$ and $Y\neq 0$.}
%     \label{alg:tangentRuleProjectiveMontgomery}
%     \begin{algorithmic}
%     \Procedure{$TangentRule$}{$[X:Y:Z]$}
%         \State \todo{Write Alg.}
%         \State \Return $([X':Y':Z'])$
%     \EndProcedure
%     \end{algorithmic}
% \end{Algorithmus}
% 
% \begin{Algorithmus} %Use the environment only if you want to place the algorithm similar to graphics from TeX
%     \caption{Projective Chord Rule used to compute $P_1\odot_M^{Pro} P_2$ where $P_1=[X_1:Y_1:Z_1], P_2=[X_2:Y_2:Z_2]\in M_{A,B}(\mathbb{F}P^2)\setminus\{[0:1:0]\}$ and $X_1\neq X_2$.}
%     \label{alg:chordRuleProjectiveMontgomery}
%     \begin{algorithmic}
%     \Procedure{$ChordRule$}{$[X_1:Y_1:Z_1], [X_2:Y_2:Z_2]$}
%         \State \todo{Write Alg.}
%         \State \Return $([X':Y':Z'])$
%     \EndProcedure
%     \end{algorithmic}
% \end{Algorithmus}

% \section{Short Weierstrass Curves}
% As in \cite{least_authority_moonmath_2024} we will introduce elliptic curves first in their affine representation and then in their projective representation.
% Later on, we will use both representations in the voting scheme.
% 
% \subsection{Affine Short Weierstrass form}
% We can define a Short Weierstrass elliptic curve in affine representation as follows.
% 
% \begin{definition}[Affine Short Weierstrass curve]
%     Let $(\mathbb{F}, +, \cdot)$ be a finite field of order $q$ with $q>3$. 
%     Let $a,b\in \mathbb{F}$ with $4a^2 + 27b^2\neq 0$.
%     Then, the corresponding Affine Short Weierstrass Curve is the group $(E_{a,b}(\mathbb{F}), \odot_E^{Aff})$, with:
%     \begin{itemize}
%         \item The set of points $E_{a,b}(\mathbb{F})$:
%         \begin{align*}
%             E_{a,b}(\mathbb{F})=\{(x,y)\in\mathbb{F}\times\mathbb{F}|y^2=x^3+a\cdot x+b\}\cup \{\mathcal{O}\}
%         \end{align*}
%         Here, $\mathcal{O}$ is the so-called \glqq point at infinity\grqq.
%         \item The group law $\odot_E^{Aff}$:\\
%         Let $P_1, P_2\in E_{a,b}(\mathbb{F})$.
%         Then:
%         \begin{align*}
%             P_1\odot_E^{Aff} P_2 =\begin{cases}
%                 P_1, &\quad\text{if } P_2=\mathcal{O}\\
%                 P_2, &\quad\text{if } P_1=\mathcal{O}\\
%                 \mathcal{O}, &\quad\text{if } P_1=(x, y) \text{ and } P_2=(x, -y)\\
%                 TangentRule(x,y), &\quad\text{if } P_1=P_2=(x,y) \text{ and } y\neq 0\\
%                 ChordRule(x_1, y_1, x_2, y_2), &\quad\text{otherwise, with } P_1=(x_1, y_1), P_2=(x_2, y_2)
%                 %\left(\left(\frac{3x^2+a}{2y}\right)^2-2x, \left(\frac{3x^2+a}{2y}\right)\cdot\left(x-\left(\left(\frac{3x^2+a}{2y}\right)^2-2x\right)\right)-y\right), &\quad\text{if } P_1=P_2=(x,y) \text{ and } y\neq 0\\
%                 %\left(\left(\frac{y_2-y_1}{x_2-x_2}\right)^2-x_1-x_2, \left(\frac{y_2-y_1}{x_2-x_1}\right)\cdot\left(x_1-\left(\left(\frac{y_2-y_1}{x_2-x_2}\right)^2-x_1-x_2\right)\right)-y_1\right), &\quad\text{otherwise, with } P_1=(x_1, y_1), P_2=(x_2, y_2)
%             \end{cases}
%         \end{align*}
%         $TangentRule$ and $ChordRule$ are defined in \ref{alg:tangentRuleAffineWeierstrass} and \ref{alg:chordRuleAffineWeierstrass}.
%     \end{itemize}
%     The neutral element of $E_{a,b}(\mathbb{F})$ is $\mathcal{O}$.
%     For $P\in E_{a,b}(\mathbb{F})$, the inverse element $P^{-1}$ is defined as $P^{-1}:=\begin{cases}
%         \mathcal{O}, &\quad\text{if } P=\mathcal{O}\\
%         (x, -y), &\quad\text{otherwise, with } P=(x,y)
%     \end{cases}$.
% \end{definition}
% 
% \begin{Algorithmus} %Use the environment only if you want to place the algorithm similar to graphics from TeX
%     \caption{Affine Tangent Rule used to compute $P\odot_E^{Aff} P$ where $P=(x,y)\in E_{a,b}(\mathbb{F})\setminus\{\mathcal{O}\}$ and $y\neq 0$.}
%     \label{alg:tangentRuleAffineWeierstrass}
%     \begin{algorithmic}
%     \Procedure{$TangentRule$}{$x,y$}
%         \State $tmp := \frac{3x^2+a}{2y}$
%         \State $x' := tmp^2-2x$
%         \State $y':= tmp\cdot (x-x')-y$
%         \State \Return $(x', y')$
%     \EndProcedure
%     \end{algorithmic}
% \end{Algorithmus}
% 
% \begin{Algorithmus} %Use the environment only if you want to place the algorithm similar to graphics from TeX
%     \caption{Affine Chord Rule used to compute $P_1\odot_E^{Aff} P_2$ where $P_1=(x_1,y_1), P_2=(x_2, y_2)\in E_{a,b}(\mathbb{F})\setminus\{\mathcal{O}\}$ and $x_1\neq x_2$.}
%     \label{alg:chordRuleAffineWeierstrass}
%     \begin{algorithmic}
%     \Procedure{$ChordRule$}{$x_1,y_1, x_2, y_2$}
%         \State $tmp := \frac{y_2-y_1}{x_2-x_1}$
%         \State $x' := tmp^2-x_1-x_2$
%         \State $y':= tmp\cdot (x_1-x')-y_1$
%         \State \Return $(x', y')$
%     \EndProcedure
%     \end{algorithmic}
% \end{Algorithmus}
% 
% \subsection{Projective Short Weierstrass form}
% To avoid the special point at infinity $\mathcal{O}$ there exists the projective representation of elliptic curves based on projective planes.
% 
% \begin{definition}[Projective Point and Projective Plane]
%     Let $(\mathbb{F}, +, \cdot)$ be a field. 
%     We define $\mathbb{F}^*$ as the set of invertible elements in $\mathbb{F}$.
%     Let $(X, Y, Z)=\in\mathbb{F}^3$. 
%     Then, the projective point corresponding to $(X, Y, Z)$ is defined as
%     \begin{align*}
%         [X:Y:Z] :=\{(k\cdot X, k\cdot Y, k\cdot Z)|k\in\mathbb{F}^*\}\text{.}
%     \end{align*}
%     The projective plane of $\mathbb{F}$ is defined as
%     \begin{align*}
%         \mathbb{F}P^2:=\{[X:Y:Z]|(X,Y,Z)\in\mathbb{F}^3 \text{ and } (X, Y, Z)\neq (0,0,0)\}\text{.}
%     \end{align*}
% \end{definition}
% 
% With this, we can define Short Weierstrass Elliptic Curves in their projective representation as groups of projective points.
% 
% \begin{definition}[Projective Short Weierstrass Curve]
%     Let $(\mathbb{F}, +, \cdot)$ be a finite field of order $q$ with $q>3$. 
%     Let $a,b\in \mathbb{F}$ with $4a^2 + 27b^2\neq 0$. 
%     Let $\mathbb{F}P^2$ be the projective plane of $\mathbb{F}$.
%     Then, the corresponging Projective Short Weierstrass Curve is the group $(E_{a,b}(\mathbb{F}P^2), \odot_E^{Pro})$, with:
%     \begin{itemize}
%         \item The set of points $E_{a,b}(\mathbb{F}P^2)$:
%         \begin{align*}
%             E_{a,b}(\mathbb{F}P^2)=\{[X:Y:Z]\in\mathbb{F}P^2|Y^2\cdot Z=X^3+a\cdot X\cdot Z^2+b\cdot Z^3\}
%         \end{align*}
%         Here, $[0:1:0]$ corresponds to the point at infinity $\mathcal{O}$ of the affine form.
%         \item The group law $\odot_E^{Pro}$:\\
%         Let $P_1=[X_1:Y_1:Z_1], P_2=[X_2:Y_2:Z_2]\in E_{a,b}(\mathbb{F}P^2)$. 
%         Then:
%         \begin{align*}
%             P_1\odot_E^{Pro} P_2 =\begin{cases}
%                 P_1, &\quad\text{if } P_2=[0:1:0]\\
%                 P_2, &\quad\text{if } P_1=[0:1:0]\\
%                 [0:1:0], &\quad\text{if } [X_1:Y_1:Z_1]=[X_2:-Y_2:Z_2]\\
%                 TangentRule([X_1:Y_1:Z_1]), &\quad\text{if } P_1=P_2 \text{ and } Y_1\neq 0\\
%                 ChordRule([X_1:Y_1:Z_1], [X_2:Y_2:Z_2]), &\quad\text{otherwise}
%                 %\left(\left(\frac{3x^2+a}{2y}\right)^2-2x, \left(\frac{3x^2+a}{2y}\right)\cdot\left(x-\left(\left(\frac{3x^2+a}{2y}\right)^2-2x\right)\right)-y\right), &\quad\text{if } P_1=P_2=(x,y) \text{ and } y\neq 0\\
%                 %\left(\left(\frac{y_2-y_1}{x_2-x_2}\right)^2-x_1-x_2, \left(\frac{y_2-y_1}{x_2-x_1}\right)\cdot\left(x_1-\left(\left(\frac{y_2-y_1}{x_2-x_2}\right)^2-x_1-x_2\right)\right)-y_1\right), &\quad\text{otherwise, with } P_1=(x_1, y_1), P_2=(x_2, y_2)
%             \end{cases}
%         \end{align*}
%         $TangentRule$ and $ChordRule$ are defined in \ref{alg:tangentRuleProjectiveWeierstrass} and \ref{alg:chordRuleProjectiveWeierstrass}.
%     \end{itemize}
%     The neutral element of $E_{a,b}(\mathbb{F}P^2)$ is $[0:1:0]$.
%     For $P=[X:Y:Z]\in E_{a,b}(\mathbb{F}P)$, the inverse element $P^{-1}$ is defined as $P^{-1}:=[X:-Y:Z]$.
% \end{definition}
% 
% \begin{Algorithmus} %Use the environment only if you want to place the algorithm similar to graphics from TeX
%     \caption{Projective Tangent Rule used to compute $P\odot_E^{Pro} P$ where $P=[X:Y:Z]\in E_{a,b}(\mathbb{F})\setminus\{[0:1:0]\}$ and $Y\neq 0$.}
%     \label{alg:tangentRuleProjectiveWeierstrass}
%     \begin{algorithmic}
%     \Procedure{$TangentRule$}{$[X:Y:Z]$}
%         \State $W := a\cdot Z^2 + 3\cdot X^3$
%         \State $S := Y\cdot Z$
%         \State $B := X\cdot Y\cdot S$
%         \State $H := W^2-8\cdot B$
%         \State $X':= 2\cdot H\cdot S$
%         \State $Y':= W\cdot (4\cdot B-H)-8\cdot Y^2\cdot S^2$
%         \State $Z':= 8\cdot S^3$
%         \State \Return $([X':Y':Z'])$
%     \EndProcedure
%     \end{algorithmic}
% \end{Algorithmus}
% 
% \begin{Algorithmus} %Use the environment only if you want to place the algorithm similar to graphics from TeX
%     \caption{Projective Chord Rule used to compute $P_1\odot_E^{Pro} P_2$ where $P_1=[X_1:Y_1:Z_1], P_2=[X_2:Y_2:Z_2]\in E_{a,b}(\mathbb{F}P^2)\setminus\{[0:1:0]\}$ and $X_1\neq X_2$.}
%     \label{alg:chordRuleProjectiveWeierstrass}
%     \begin{algorithmic}
%     \Procedure{$ChordRule$}{$[X_1:Y_1:Z_1], [X_2:Y_2:Z_2]$}
%         \State $U_1 := Y_2\cdot Z_1$
%         \State $U_2 := Y_1\cdot Z_2$
%         \State $U := U_1 - U_2$
%         \State $V_1 := X_2\cdot Z_1$
%         \State $V_2 := X_1\cdot Z_2$
%         \State $V := V_1 - V_2$
%         \State $W := Z_1\cdot Z_2$
%         \State $A := U^2\cdot W - V^3-2\cdot V^2\cdot V_2$
%         \State $X' := V\cdot A$
%         \State $Y' := U\cdot (V^2\cdot V_2 - A)-V^3\cdot U_2$
%         \State $Z' := V^3\cdot W$
%         \State \Return $([X':Y':Z'])$
%     \EndProcedure
%     \end{algorithmic}
% \end{Algorithmus}

\section{Exponentiation}
We use a Montgomery curve $M_{A,B}(\mathbb{F}P^2)$ as group in \gls{eeg} encryption because computing exponentiations of group elements is comparatively efficient.
Group exponentiation is defined as follows:
\begin{align*}
    \forall P=[X:Y:Z]\in M_{A,B}(\mathbb{F}P^2)\forall n\in\mathbb{N}:\quad P^n:=\underbrace{P\odot_{M}^{Pro}\dots\odot_{M}^{Pro}P}_{n \text{ times}}
\end{align*}
However, computing $P^n=[X_n:Y_n:Z_n]$ by iteratively using the group law $\odot_{M}^{Pro}$ $n$ times is computationally expensive.
Therfore, we utilize the Montgomery Ladder algorithm \cite{costello_montgomery_2018}.
This algorithm only requires $\mathcal{O}(\log_2(n))$ many iterations to compute the coordinates $X_n$ and $Z_n$ of $P^n$ from the coordinates $X$ and $Z$ of $P$.
Afterwards, we use the Okeya-Sakurai $Y$-Recovery algorithm \cite{goos_efficient_2001} to reconstruct $Y_n$ from the last two iterations of the Montgomery Ladder.

\subsection{Pseudooperations}
For the Montgomery Ladder algorithm we work only with the $X$ and $Z$ coordinates of points $P=[X:Y:Z]\in M_{A,B}(\mathbb{F}P^2)$.
For this purpose, we define a quotient map $x$ to map projective points to pairs of their $X$ and $Z$ coordinates.
Specially treated is the point at infinity $[0:1:0]$.
\begin{definition}[Quotient Map of Montgomery Curves]
    Let $M_{A,B}(\mathbb{F}P^2)$ be a Montgomery curve.
    Then, we define te quotient map $x$ as follows:
    \begin{align*}
        &x:M_{A,B}(\mathbb{F}P^2)\to \{[X:Z]|X,Z\in\mathbb{F}\},\\
        &x([X:Y:Z]) :=\begin{cases}
            [X:Z], \quad\text{if } Z\neq 0\\
            [1:0], \quad\text{otherwise ($\rightsquigarrow [X:Y:Z]=[0:1:0]$)}
        \end{cases}
    \end{align*}
    We call $x(P)$ the Pseudopoint of $P\in M_{A,B}(\mathbb{F}P^2)$.
\end{definition}

Using this map, we want to be able to apply the group law to Pseudopoints. 
In other words, we want to compute $x(P\odot_M^{Pro}Q)$ from $x(P)$ and $x(Q)$.
Unfortunately, as mentioned in chapter 3.1 of \cite{costello_montgomery_2018} this is not possible, since $x(P)$ and $x(Q)$ doesn't uniquely determine $x(P\odot_M^{Pro} Q)$.
However, it is possible to determine $x(P\odot_M^{Pro} Q)$ based on $x(P), x(Q)$ and $x(P\odot_M^{Pro}Q^{-1})$.
This is useful for the Montgomery Ladder later on, since $P\odot_M^{Pro}Q$ will be a fixed, known point.

Based on this, we can define two pseudo-operations $xAdd$ and $xDbl$ which mostly parallel the Chord Rule and Tangent Rule of the group law $\odot_M^{Aff}$.
We define these operations analog to the approach presented in chapter 3 of \cite{costello_montgomery_2018}.

\begin{definition}[Pseudooperation $xAdd$]
    Let $M_{A,B}(\mathbb{F}P^2)$ be a Montgomery curve.
    Let $P, Q\in M_{A,B}(\mathbb{F}P^2)$ with $P\neq Q$.
    We define:
    \begin{align*}
        xAdd(x(P), x(Q), x(P\odot_M^{Pro}Q^{-1})) := x(P\odot_M^{Pro} Q)
    \end{align*}
    We can compute $xAdd$ as presented in \cref{alg:xAdd}.
\end{definition}

\begin{Algorithmus} %Use the environment only if you want to place the algorithm similar to graphics from TeX
    \caption{Pseudooperation $xAdd$ used to compute $x(P\odot_M^{Pro} Q)$ from $x(P), x(Q)$ and $x(P\odot_M^{Pro}Q^{-1})$ with $P,Q\in M_{A,B}(\mathbb{F}P^2)$ and $P\neq Q$.}
    \label{alg:xAdd}
    \begin{algorithmic}
    \Procedure{$xAdd$}{$x(P), x(Q), (P\odot_M^{Pro}Q^{-1})$}
        \State $[X_P:Z_P]:=x(P);\: [X_Q:Z_Q]:=x(Q);\: [X_{\odot^{-1}}, Z_{\odot^{-1}}]:=x(P\odot_m^{Pro}Q^{-1})$
        \State $X_\odot:= Z_{\odot^{-1}}((X_P-Z_P)(X_Q+Z_Q)+(X_P+Z_P)(X_Q-Z_Q))^2$
        \State $Z_\odot:= X_{\odot^{-1}}((X_P-Z_P)(X_Q+Z_Q)-(X_P+Z_P)(X_Q-Z_Q))^2$
        \State \Return $[X_\odot:Z_\odot]$
    \EndProcedure
    \end{algorithmic}
\end{Algorithmus}

\begin{definition}[Pseudooperation $xDbl$]
    Let $M_{A,B}(\mathbb{F}P^2)$ be a Montgomery curve.
    Let $P\in M_{A,B}(\mathbb{F}P^2)$.
    We define:
    \begin{align*}
        xDbl(x(P)) := x(P\odot_M^{Pro} P) = x(P^2)
    \end{align*}
    We can compute $xDbl$ as presented in \cref{alg:xDbl}.
\end{definition}

\begin{Algorithmus} %Use the environment only if you want to place the algorithm similar to graphics from TeX
    \caption{Pseudooperation $xDbl$ used to compute $x(P\odot_M^{Pro} P)=x(P^2)$ from $x(P)$ with $P\in M_{A,B}(\mathbb{F}P^2)$.}
    \label{alg:xDbl}
    \begin{algorithmic}
    \Procedure{$xDbl$}{$x(P)$}
        \State $[X_P:Z_P]:=x(P)$
        \State $X_\odot:= (X_P+Z_P)^2(X_P-Z_P)^2$
        \State $Z_\odot:= \left(4X_PZ_P\right)\left(\left(X_P-Z_P\right)^2+\left(\frac{A+2}{4}\left(4X_PZ_P\right)\right)\right)$
        \State \Return $[X_\odot:Z_\odot]$
    \EndProcedure
    \end{algorithmic}
\end{Algorithmus}

\subsection{The Montgomery Ladder}
The Idea behind the Montgomery Ladder is similar to that of Fast Exponentiation.
Analogously to Fast Exponentiation, we can compute exponentiations $P^n$ for some $P\in M_{A,B}(\mathbb{F}P^2)$ and $n\in\mathbb{N}$.
However, we can also employ $xAdd$ and $xDbl$ to more efficiently determine $x(P^n)$ based on $x(P)$. 
The Montgomery Ladder Algorithm computes $x(P^n)$ along with $x(P^{n+1})$ which we need later on to recover the $Y$ coordinate of $P^n$.

\begin{definition}[Montgomery Ladder]
    Let $M_{A,B}(\mathbb{F}P^2)$ be a Montgomery curve.
    Let $P\in M_{A,B}(\mathbb{F}P^2)$.
    Let $n\in\mathbb{N}$, where $n=\sum\limits_{i=0}^{l-1}n_i2^i$ is a bit string of fixed length $l$.
    Then, we can compute $(x(P^n), x(P^{n+1}))$ with $MontgomeryLadder$ defined in \cref{alg:montgomeryLadder}.
\end{definition}

\begin{Algorithmus} %Use the environment only if you want to place the algorithm similar to graphics from TeX
    \caption{Montgomery Ladder algorithm used to compute $x(P^n)$ and $x(P^{n+1})$ from $x(P)$ with $P\in M_{A,B}(\mathbb{F}P^2)$ and $n\in\mathbb{N}$.
        Here, $n=\sum\limits_{i=0}^{l-1}n_i2^i$ is a bit string of fixed length $l$.
    }
    \label{alg:montgomeryLadder}
    \begin{algorithmic}
    \Procedure{$MontgomeryLadder$}{$x(P), (n_0,\dots, n_{l-1})$}
        \State $(R_0, R_1):=(x([0:1:0]), x(P))$ 
        \State \Comment \parbox[t]{0.7\linewidth}{%
            Initialize $R_0$ with point at infinity and $R_0$ with $x(P)$. This is different compared to the standard version of the Montgomery Ladder.%
        }
        \For{$i=l-1$ \textbf{down to} $0$}
            \If{$n_i=0$}
                \State $(R_0, R_1):=(xDbl(R_0), xAdd(R_1, R_0, x(P)))$ 
                \State \Comment \parbox[t]{0.7\linewidth}{%
                    $R_1=x(P^{k+1})$ and $R_0=x(P^k)$ for some $k\in\mathbb{N}$ at all times.\\
                    Therefore, $x(P^{k+1}\odot_{M}^{Pro}(P^k)^{-1})=x(P)$.
                }
            \Else
                \State $(R_0, R_1):=(xAdd(R_1, R_0, x(P)), xDbl(R_1))$
            \EndIf
        \EndFor
        \State \Return $(R_0, R_1)$
    \EndProcedure
    \end{algorithmic}
\end{Algorithmus}

Note that this is not the standard version of the Montgomery Ladder but slightly changed and originally intendended for constant-time execution \cite{costello_montgomery_2018}.
In the standard version of the Montgomery Ladder algorithm, we require the most significant bit of the exponent $n_{l-1}$ to be $1$. 
However, this means that the length of the bitstring representing $n$ changes based on the exponent $n$.
When we implement the Montgomery Ladder in a ciruit for our voting scheme this is problematic.
Since the exponent is not known during the compilation time of the circuit, we need to work with a exponents of a fixed length $l$ where the most significatn bit might be zero.
Therefore, this version of the Montgomery Ladder is better suited for our purposes.

\subsection{$Y$-Recovery}
After performing the Montgomery Ladder algorithm on $x(P)$ and $n\in\mathbb{N}$ with $P\in M_{A,B}(\mathbb{F}P^2)$, we have $[X_n:Z_n]:=x(P^n)$ and $[X_{n+1}:Z_{n+1}]:=x(P^{n+1})$.
To reconstruct the $Y$ coordinate $Y_n$ of $P^{n}=[X_n:Y_n:Z_n]$, we use a recovery algorithm based on $P, x(P^n)$ and $x(P^{n+1})$ introduced by Okeya and Sakurai in \cite{goos_efficient_2001}.

\begin{definition}[Okeya-Sakurai $Y$-Recovery]
    Let $M_{A,B}(\mathbb{F}P^2)$ be a Montgomery curve.
    Let $P=[X_P:Y_P:Z_P], Q, P\odot_{M}^{Pro}Q\in M_{A,B}(\mathbb{F}P^2)\setminus\{[0:1:0]\}$ with $Y_P\neq 0$.
    Then, we can reconstruct the complete $Q=Q_{rec}$, including the $Y$ coordinate from $(x_P, y_P):=I^{-1}(P)=\left(\frac{X_P}{Z_P}, \frac{Y_P}{Z_P}\right)$, $x(Q)=[X_Q:Z_Q]$ and $x(P\odot_M^{Pro} Q)=[X_\odot:Z_\odot]$ using \cref{alg:okSakYRecovery}.
\end{definition}

\begin{Algorithmus} %Use the environment only if you want to place the algorithm similar to graphics from TeX
    \caption{Okeya-Sakurai $Y$-Recovery used to recover $Q$ with $Y$ coordinate from $P$, $x(Q)$ and $x(P\odot_M^{Pro} Q)$.
    Here, $P=[X_P:Y_P:Z_P],P\odot_{M}^{Pro}Q\in M_{A,B}(\mathbb{F}P^2)\setminus\{[0:1:0]\}$, with $Y_P\neq 0$.}
    \label{alg:okSakYRecovery}
    \begin{algorithmic}
    \Procedure{$OkSakYRecovery$}{$P, x(Q), x(P\odot_M^{Pro}Q)$}
        \State $[X_P:Y_P:Z_P]:=P;\: [X_Q:Z_Q]:=x(Q);\: [X_\odot:Z_\odot]:=x(P\odot_M^{Pro}Q)$
        \State $(x_P, y_P):=\left(\frac{X_P}{Z_P}, \frac{Y_P}{Z_P}\right)$
        \State $X_{Q_{rec}}:=2By_PZ_QZ_\odot X_Q$
        \State $Y_{Q_{rec}}:=Z_\odot((X_Q+x_PZ_Q+2AZ_Q)(X_Qx_P+Z_Q)-2AZ_Q^2)-(X_Q-x_PZ_Q)^2X_\odot$
        \State $Z_{Q_{rec}}:=2By_PZ_Q^2Z_\odot$
        \State $Q_{rec} := [X_{Q_{rec}}:Y_{Q_{rec}}:Z_{Q_{rec}}]$
        \State \Return $Q_{rec}$
    \EndProcedure
    \end{algorithmic}
\end{Algorithmus}

Now we can recover the $Y$ coordinate of a point $Q$ from its pseudopoint $x(Q)$, $P$ and $x(P\odot_M^{Pro} Q)$ if the preconditions for Okeya-Sakurai $Y$-Recovery are fulfilled.
However, we only want to employ the Okeya-Sakurai $Y$-Recovery for group exponentiations.
Here, we want to compute $P^n$ for some $P\in M_{A,B}^{Pro}$ and $n\in \mathbb{N}$.
The Montgomery Ladder produces the pseudopoints $x(P^n)$ and $x(P^{n+1})$.
In order to recover the correct $Y$-Coordinate of $P^n$, we can use the Okeya-Sakurai $Y$-Recovery with $P:=P$, $Q:=P^n$ and $P\odot_M^{Pro} Q:=P^{n+1}$ but need to treat some edge cases seperately:
\begin{itemize}
    \item First, we consider the case that $P=[0:1:0]$. 
    Here, we can conclude $P^{n+1}=P^n=[0:1:0]$.
    Note that the $OkSakYRecovery$ returns the correct result $P^n=Q_{rec}=[X_{Q_{rec}}:Y_{Q_{rec}}:Z_{Q_{rec}}]=[0:Y_{Q_{rec}}:0]=[0:1:0]$ as long as $Y_{Q_{rec}}\neq 0$.
    We introduce a normalization function $NormProjective$ which covers the edge case $Y_{Q_{rec}}=0$ and does not change any other projective coordinates.
    \begin{align*}
        NormProjective:\{\mathbb{F}^3\}\to\{\mathbb{F}^3\}, [X:Y:Z]\to\begin{cases}
            \left[\frac{X}{Z}:\frac{Y}{Z}:1\right], &\:\text{if } Z\neq 0\\
            [0:1:0], &\:\text{otherwise}
        \end{cases}
    \end{align*}
    If we always perform this operation on the output of $OkSakYRecovery$, we don't need to check for $P^n=[0:1:0]$.
    \item Next, we know that $P^{n+1}=[0:1:0]\:\Leftrightarrow\:P^n=P^{-1}$.
    In this case, the correct output would be $P^n=Q_{rec}=P^{-1}=[X_P:-Y_P:Z_P]$.
\end{itemize}
% However, we only want to employ the Okeya-Sakurai $Y$ recovery for group exponentiations.
% Therefore, we have the special case, where $P\in M_{A,B}^{Pro}(\mathbb{F}P^2)$, $Q=P^n$ and $P\odot_M^{Pro}Q=P^{n+1}$ for some $n\in\mathbb{N}$.
% In this case, the conditions $P=[X_P:Y_P:Z_P], Q, P\odot_M^{Pro}Q\in M_{A,B}^{Pro}(\mathbb{F}P^2)\setminus\{[0:1:0]\}$ and $Y_P\neq 0$ simplify to 

We mostly follow the implementation in \cite{huber_hicolasnuberballotsnarks_2024} for the $Y$-Recovery on the output of the Montgomery Ladder.

\begin{definition}[Montgomery Ladder $Y$-Recovery]
    Let $M_{A,B}(\mathbb{F}P^2)$ be a Montgomery curve.
    Let $P\in M_{A,B}(\mathbb{F}P^2)$, $n\in\mathbb{N}$.
    Then, we can reconstruct the complete $P^n=Q_{rec}$, including the $Y$ coordinate from $P, x(P^n)$ and $x(P^{n+1})$ using \cref{alg:montgomeryYRecovery}.
\end{definition}

\begin{Algorithmus} %Use the environment only if you want to place the algorithm similar to graphics from TeX
    \caption{$Y$-Recovery used to recover $P^n$ with $Y$ coordinate from $P$, $x(P^n)$ and $x(P^{n+1})$ for some $n\in\mathbb{N}$.
    Here, $P\in M_{A,B}(\mathbb{F}P^2)$.}
    \label{alg:montgomeryYRecovery}
    \begin{algorithmic}
    \Procedure{$LadderYRecovery$}{$P, x(P^n), x(P^{n+1})$}
        \State $[X_P:Y_P:Z_P]:=P;\: [X_{P^n}:Z_{P^n}]:=x(P^n);\: [X_{P^{n+1}}:Z_{P^{n+1}}]:=x(P^{n+1})$
        \State $P_{rec}^n:=[0:1:0]$
        \If{$Z_{P^{n+1}}=0$} \Comment \parbox[t]{0.7\linewidth}{%
                $\rightsquigarrow P^{n+1}=[0:1:0]$ and therefore $P^n=P^{-1}$.
            }
            \State $P_{rec}^n:=[X_P:-Y_P:Z_P]$
        \Else
            \State $P_{rec}^n:=OkSakYRecovery(P, x(P^n), x(P^{n+1}))$
        \EndIf
        \State $P_{rec}^n:=NormProjective(P_{rec}^n)$
        \State \Return $P_{rec}$
    \EndProcedure
    \end{algorithmic}
\end{Algorithmus}

\subsection{Full Exponentiation}
Now, we can combine the Montgomery Ladder with the $Y$-Recovery to compute full exponentiations $P^n$ for some $P\in M_{A,B}(\mathbb{F}P^2)$ and $n\in\mathbb{N}$.

\begin{definition}[Montgomery Exponentiation]
    Let $M_{A,B}(\mathbb{F}P^2)$ be a Montgomery curve.
    Let $P\in M_{A,B}(\mathbb{F}P^2)$.
    Let $n\in\mathbb{N}$, where $n=\sum\limits_{i=0}^{l-1}n_i2^i$ is a bit string of fixed length $l$.
    Then, we can compute $P^n$ with $Exponentiation$ as described in \cref{alg:montgomeryExponentiation}.
\end{definition}

\begin{Algorithmus} %Use the environment only if you want to place the algorithm similar to graphics from TeX
    \caption{Exponentiation on Montgomery curves. Used to compute $P^n$ where $P\in M_{A,B}^{Pro}$ and $n\in\mathbb{N}$. 
        Here, $n=\sum\limits_{i=0}^{l-1}n_i2^i$ is a bit string of fixed length $l$.
    }
    \label{alg:montgomeryExponentiation}
    \begin{algorithmic}
    \Procedure{$Exponentiation$}{$P, (n_0,\dots,n_{l-1})$}
        \State $(x(P^n), x(P^{n+1})):= MontgomeryLadder(x(P), (n_0,\dots,n_{l-1}))$
        \State $P^n:=LadderYRecovery(P, x(P^n), x(P^{n+1}))$
        \State \Return $P^n$
    \EndProcedure
    \end{algorithmic}
\end{Algorithmus}